\section{Worked cases}
\label{sec:reference}

\subsection{The cases}
Four operating cases are worked through below: the 1\,GW symmetrical monopole,
a double symmetrical monopole variant, the 3\,GW bipole, and that bipole
holding half its transfer on one pole. Station losses are 1.2\,\% per station
and the current margin 5\,\%. All conductor data is the DC resistance at
20\,$^{\circ}$C from the ACSR tables and the flow is evaluated at a conductor
temperature of 75\,$^{\circ}$C.

\begin{table}[t]
  \centering
  \small
  \renewcommand{\arraystretch}{1.25}
  \begin{tabularx}{\linewidth}{@{}X l r r r r@{}}
    \toprule
    \textbf{Given} & \textbf{Unit} &
    \textbf{SMP} & \textbf{DSMP} & \textbf{BIP} &
    \textbf{BIP 1 pole} \\
    \midrule
    AC power at receiving end     & MW                & 1{,}000 & 1{,}500 & 3{,}000 & 1{,}500 \\
    DC voltage pole-to-ground $U_{dpe,rec}$ & kV                & 400     & 400     & 525     & 525     \\
    DC line length                & mi                & 160     & 200     & 200     & 200     \\
    HV conductor                  & ---               & Bluebird & Chukar & Bluebird & Bluebird \\
    HV subconductors per pole     & ---               & 1       & 2       & 2       & 2       \\
    Specific HV resistance        & $\Omega$/1000\,ft & 0.00801 & 0.00970 & 0.00801 & 0.00801 \\
    DMR conductor                 & ---               & n/a     & n/a     & Chukar  & Chukar  \\
    DMR subconductors             & ---               & n/a     & n/a     & 2       & 1       \\
    Specific DMR resistance       & $\Omega$/1000\,ft & n/a     & n/a     & 0.00970 & 0.00970 \\
    Reference temperature $T_{ref}$ & $^{\circ}$C     & 20      & 20      & 20      & 20      \\
    Conductor temperature $T$     & $^{\circ}$C       & 75      & 75      & 75      & 75      \\
    Zero-resistance temperature $T_{0}$ & $^{\circ}$C & 228     & 228     & 228     & 228     \\
    Losses per converter station  & \%                & 1.2     & 1.2     & 1.2     & 1.2     \\
    Margin for DC current         & \%                & 5       & 5       & 5       & 5       \\
    \bottomrule
  \end{tabularx}
  \caption{Given quantities for the four cases. BIP\,1\,pole is the 3\,GW
           bipole with one pole out of service, on the same plant as the BIP
           case and with its own post-contingency transfer.}
  \label{tab:refinputs}
\end{table}

\subsection{Temperature correction and line resistance}
The correction factor is common to all four cases,
Equation~\eqref{eq:ktvalue}:
\[
  k_{T} = \frac{228+75}{228+20} = 1.2218 .
\]
Applying Equation~\eqref{eq:resistance} to the bipole, whose HV conductor and
length the outage case shares:
\[
  R_{d,hv} = \frac{0.00801\cdot 5.28}{2}\cdot 200 \cdot 1.2218
         = 4.229 \cdot 1.2218 = 5.167\ \Omega,
\]
and to the single-subconductor DMR that carries the return in the outage case:
\[
  R_{d,dmr} = \frac{0.00970\cdot 5.28}{1}\cdot 200 \cdot 1.2218
          = 10.243 \cdot 1.2218 = 12.515\ \Omega .
\]
So $R_{loop} = 10.335\,\Omega$ balanced and $17.682\,\Omega$ with the return
through the DMR. At the 20\,$^{\circ}$C data temperature those figures would
have been $8.459\,\Omega$ and $14.472\,\Omega$: the correction alone adds
22\,\% to every resistive loss reported below.

\subsection{Balanced bipole, step by step}
Following Section~\ref{sec:solution} with $k=2$:
\begin{align*}
  P_{dc,inv} &= \frac{3000}{1-0.012} = 3036.4\ \mathrm{MW}
              && \text{Eq.~\eqref{eq:step1}} \\
  I_{d}      &= \frac{525 - \sqrt{525^{2} - 10.335\cdot 3036.4}}{10.335}
              = \frac{525 - 494.21}{10.335} = 2.979\ \mathrm{kA}
              && \text{Eq.~\eqref{eq:current2}} \\
  P_{dc,rec} &= 2\cdot 525 \cdot 2.979 = 3128.2\ \mathrm{MW}
              && \text{Eq.~\eqref{eq:c1}} \\
  P_{loss,line} &= 10.335 \cdot 2.979^{2} = 91.7\ \mathrm{MW}
              && \text{Eq.~\eqref{eq:c2}} \\
  U_{dpe,inv}  &= 525 - 2.979\cdot 5.167 = 509.6\ \mathrm{kV}
              && \text{Eq.~\eqref{eq:vinv}} \\
  P_{ac,rec} &= \frac{3128.2}{1-0.012} = 3166.2\ \mathrm{MW}
              && \text{Eq.~\eqref{eq:c4}} \\
  \eta       &= \frac{3000}{3166.2} = 94.75\,\%
              && \text{Eq.~\eqref{eq:c6}}
\end{align*}
The loss balance of Equation~\eqref{eq:closure} closes:
$38.0 + 91.7 + 36.4 = 166.2\ \mathrm{MW} = P_{ac,rec}-P_{ac,inv}$.

\subsection{One pole out, step by step}
The same plant with $k=1$, the DMR in the loop, and its own 1500\,MW target:
\begin{align*}
  P_{dc,inv} &= \frac{1500}{1-0.012} = 1518.2\ \mathrm{MW} \\
  I_{d}      &= \frac{525 - \sqrt{525^{2} - 4\cdot 17.682\cdot 1518.2}}
                     {2\cdot 17.682}
              = \frac{525 - 410.16}{35.36} = 3.247\ \mathrm{kA}
              && \text{Eq.~\eqref{eq:current}} \\
  P_{dc,rec} &= 525 \cdot 3.247 = 1704.6\ \mathrm{MW} \\
  P_{loss,line} &= 5.167\cdot 3.247^{2} + 12.515\cdot 3.247^{2}
                 = 54.5 + 131.9 = 186.4\ \mathrm{MW}
              && \text{Eq.~\eqref{eq:losssplit}}
\end{align*}
Half the power, and yet a higher current and twice the line loss of the
balanced bipole. This is the case that sizes the conductor.

\subsection{Results}
Table~\ref{tab:refresults} collects the four cases, each at the
75\,$^{\circ}$C conductor temperature.

\begin{table}[t]
  \centering
  \small
  \renewcommand{\arraystretch}{1.25}
  \begin{tabularx}{\linewidth}{@{}X l r r r r@{}}
    \toprule
    \textbf{Result} & \textbf{Unit} &
    \textbf{SMP} & \textbf{DSMP} & \textbf{BIP} &
    \textbf{BIP 1 pole} \\
    \midrule
    Temperature correction factor    & ---      & 1.2218  & 1.2218  & 1.2218  & 1.2218  \\
    HV resistance per pole           & $\Omega$ & 8.268   & 6.257   & 5.167   & 5.167   \\
    DMR resistance                   & $\Omega$ & ---     & ---     & ---     & 12.515  \\
    Loop resistance                  & $\Omega$ & 16.535  & 12.515  & 10.335  & 17.682  \\
    DC power at inverter terminals   & MW       & 1{,}012.1 & 1{,}518.2 & 3{,}036.4 & 1{,}518.2 \\
    DC current, operating            & A        & 1{,}300 & 1{,}958 & 2{,}979 & 3{,}247 \\
    DC current, rated (5\,\% margin) & A        & 1{,}365 & 2{,}056 & 3{,}128 & 3{,}409 \\
    DC power at rectifier terminals  & MW       & 1{,}040.1 & 1{,}566.2 & 3{,}128.2 & 1{,}704.6 \\
    Line losses, HV conductors       & MW       & 27.9    & 48.0    & 91.7    & 54.5    \\
    Line losses, DMR                 & MW       & ---     & ---     & 0.0     & 131.9   \\
    Total DC line losses             & MW       & 27.9    & 48.0    & 91.7    & 186.4   \\
    DC voltage at inverter $U_{dpe,inv}$ & kV    & 389.3   & 387.7   & 509.6   & 508.2   \\
    DC voltage at inverter $U_{dpm,inv}$ & kV       & ---     & ---     & ---     & 467.6   \\
    AC power at sending end          & MW       & 1{,}052.7 & 1{,}585.2 & 3{,}166.2 & 1{,}725.3 \\
    Total losses, AC to AC           & MW       & 52.7    & 85.2    & 166.2   & 225.3   \\
    Overall efficiency               & \%       & 94.99   & 94.62   & 94.75   & 86.94   \\
    Maximum transferable power       & MW       & 9{,}676 & 12{,}785 & 26{,}671 & 3{,}897 \\
    \bottomrule
  \end{tabularx}
  \caption{Results at a conductor temperature of 75\,$^{\circ}$C.}
  \label{tab:refresults}
\end{table}

\subsection{Reading the results}
Four points are worth drawing out.

\textbf{The outage state sets the current rating.} At 3247\,A the bipole
holding half its transfer on one pole draws 9\,\% more current than it does at
full power on two, and with the margin applied it asks for 3409\,A against
3128\,A in normal operation. A conductor or valve sized on balanced operation
alone would be undersized for the contingency it exists to survive.

\textbf{The DMR dominates the outage loss.} 132\,MW of the 186\,MW lost in the
one-pole-out case is in the DMR alone, because a single Chukar subconductor is
2.4 times the resistance of the two-conductor Bluebird pole bundle it stands in
for. Doubling the DMR to two subconductors brings $R_{d,dmr}$ to
6.257\,$\Omega$, the current to 3101\,A and the line loss to 110\,MW --- a
76\,MW improvement for one extra conductor, and the first thing to settle in
the DMR design.

\textbf{Temperature is not a detail.} Every resistance and every line loss in
Table~\ref{tab:refresults} is 22\,\% higher than the conductor table suggests,
purely because the line is evaluated hot. Carried through to the bipole
efficiency the correction is worth about half a percentage point --- the same
order as the difference between the conductor options being compared, so a
comparison made at data temperature would not be a fair one.

\textbf{Where the power is defined changes the current.} Taken as the power
\emph{leaving} the rectifier over a line at data temperature, 3000\,MW on the
bipole corresponds to $P/(2\cdot U_{dpe,rec}) = 2857\,\mathrm{A}$. Required to
\emph{arrive} at the receiving AC terminals through a hot line, the same
3000\,MW needs 2979\,A, because the current now has to cover the line and
converter losses as well. The second figure is the one to size equipment
against when the transfer is guaranteed at the point of delivery.
